\documentclass[11pt,a4paper]{article}
\usepackage[paperwidth=210mm,paperheight=500mm,margin=1.5cm]{geometry}
\usepackage{graphicx}
\usepackage{float}
\usepackage{caption}
\usepackage{subcaption}
\usepackage{amsmath}
\usepackage[hidelinks]{hyperref}
\usepackage{xcolor}

% No page numbers
\pagestyle{empty}
\makeatletter
\let\ps@plain\ps@empty
\makeatother

\setlength{\parindent}{0pt}
\setlength{\parskip}{0.5em}
\captionsetup{font=small,labelfont=bf}

\newcommand{\plotdir}{../results}

\newcommand{\resultpair}[5]{%
\begin{figure}[H]
    \centering
    \begin{subfigure}[t]{0.48\linewidth}
        \centering
        \includegraphics[width=\linewidth]{\plotdir/#1}
        \caption{#2}
    \end{subfigure}
    \hfill
    \begin{subfigure}[t]{0.48\linewidth}
        \centering
        \includegraphics[width=\linewidth]{\plotdir/#3}
        \caption{#4}
    \end{subfigure}
    \caption{#5}
\end{figure}
}

\title{Diffusion Exercise: One-Dimensional Sample Generation}
\author{}
\date{}

\begin{document}

\maketitle

This report shows the final plots from the one-dimensional diffusion model.

The target data is a mixture of two Gaussians with peaks near $-4$ and $4$. The model is a small MLP. It takes the noisy sample and the time step as input, and it predicts the noise that was added. To generate new samples, I start from Gaussian noise and run the learned reverse diffusion process.

\section*{Training Curve}

\begin{figure}[H]
    \centering
    \includegraphics[width=0.72\linewidth]{\plotdir/loss_curve.png}
    \caption{Training and validation loss for predicting the added noise. The loss decreases quickly and then stays roughly constant, so 50 epochs is enough for this small example.}
    \label{fig:loss-curve}
\end{figure}


\section*{Sampled Distributions During Training}

\begin{figure}[H]
    \centering

    \begin{subfigure}{0.48\linewidth}
        \centering
        \includegraphics[width=\linewidth]{\plotdir/samples_epoch_0010.png}
        \caption{10 epochs}
    \end{subfigure}
    \begin{subfigure}{0.48\linewidth}
        \centering
        \includegraphics[width=\linewidth]{\plotdir/samples_epoch_0020.png}
        \caption{20 epochs}
    \end{subfigure}

    \begin{subfigure}{0.48\linewidth}
        \centering
        \includegraphics[width=\linewidth]{\plotdir/samples_epoch_0030.png}
        \caption{30 epochs}
    \end{subfigure}
    \begin{subfigure}{0.48\linewidth}
        \centering
        \includegraphics[width=\linewidth]{\plotdir/samples_epoch_0040.png}
        \caption{40 epochs}
    \end{subfigure}

    \begin{subfigure}{0.48\linewidth}
        \centering
        \includegraphics[width=\linewidth]{\plotdir/samples_epoch_0050.png}
        \caption{50 epochs}
    \end{subfigure}

    \caption{Generated samples during training. The plots are ordered left-to-right and top-to-bottom. The model first learns a rough two-peak shape. By 40--50 epochs, the samples are close to the target distribution.}
    \label{fig:samples-training}
\end{figure}


\section*{Reverse Diffusion Process}

\begin{figure}[H]
    \centering
    \includegraphics[width=0.82\linewidth]{\plotdir/diffusion_stages.png}
    \caption{Sample distributions during reverse diffusion. The samples start close to Gaussian noise. During the reverse process, they move into the two modes near $-4$ and $4$.}
    \label{fig:diffusion-stages}
\end{figure}

\end{document}
